\documentclass[12pt]{article}
\usepackage{amsmath,amsfonts,verbatim,fullpage,amsthm,graphicx}

\newtheorem*{theorem}{Theorem}
\newtheorem*{fact}{Fact}
\newtheorem*{goal}{Goal}
\newtheorem*{idea}{Idea}
\newtheorem*{defn}{Definition}
\newtheorem*{ex}{Example}

%Style section
%\setlength{\textheight}{10in}
%\setlength{\textwidth}{7in}
%\hoffset=-0.75in
 \pagestyle{plain}
% \setlength{\topmargin}{0in}
% \setlength{\headsep}{0in}
% \setlength{\oddsidemargin}{0in}
% \setlength{\evensidemargin}{0in}

\begin{document}


\pagestyle{empty} %use this to get rid of page numbers

\noindent
{Math 166 \hskip 1.5 truein  Names:\ \hrulefill}

\noindent
%\vskip 0.3truein
%\bigskip
%\noindent
{Volume via cross section}

%\noindent \hskip 3 truein  Section Number:\ \hrulefill
\medskip

\begin{goal} \rm
Use integration to find the volume of various solids.
\end{goal}
%\smallskip


%In the following questions, first graph the function(s). Determine whether to integrate with respect to $x$ or $y$. State whether the cross-section is a disk or an annulus. Then set up the integral(s). If you have time, evaluate the integral(s) to find the volume.

\begin{enumerate}
%\setcounter{enumi}{-1}

\item Set up three different integrals to find the volume of concrete in the ramp drawn below (where the dimensions are given in meters). Draw a picture to correspond to each integral. Set up your axes and be sure to draw an arbitrary cross-section (at some fixed $x$, $y$, or $z$). In one of your pictures, the cross-section should be a triangle, and in the others, a rectangle. Evaluate the integrals to see that you obtain the same answer. \\ Hint: When the cross-section is a rectangle in this example, its area will depend on the slope of the hypotenuse of the triangle.

\includegraphics[scale=.35]{concrete_ramp} 
\vspace{1in}
\\
\includegraphics[scale=.35]{concrete_ramp}
\vspace{1in}
\\
\includegraphics[scale=.35]{concrete_ramp}

%\pagebreak

%\item Calculate the volume of a cube with side length 2 by integrating. Observe that this matches with the formula you know! \\\textit{Hint: Let the cube have bottom corners $(0,0,0)$, $(2,0,0)$, $(2,2,0)$, and $(0,2,0)$.} 

\pagebreak

\item Draw the solid whose base is the triangle enclosed by the line $x+y=3$, the $y$-axis, and the $x$-axis, and whose cross-sections perpendicular to the y-axis are squares. Then use an integral to find its volume. 

\newpage

\item 
\begin{itemize}
\item Object $A$ has a base bounded by the circle $x^2 + y^2 = 4$, and its cross-sections perpendicular to the $y$-axis are squares.  
\item Object $B$ has a base bounded by the function $y = 4 - x^2$ and the $x$-axis, and its cross-sections perpendicular to the $x$-axis are rectangles with height $4$.
\end{itemize}
Explain how we know the volumes of objects $A$ and $B$ are equal without computing the volumes.

\begin{comment}
\item Draw the region in the first quadrant ($y\geq 0, x\geq 0$) bounded by $y=x^2$, the $y$-axis, and the line $y=5$. Find the volume of the solid whose base is this region and whose cross-sections perpendicular to the $y$-axis are rectangles whose height is one half the width.





\item Draw the region in the first quadrant bounded by $y=4-x^2$, the $y$-axis, and the $x$-axis. 
\begin{enumerate}
\item Then sketch the solid obtained by revolving this region about the $x$-axis. Use an integral (in the variable $x$) to find the volume of this solid.

\vspace{3.5in}
\item Now sketch the solid obtained by revolving this region about the $y$-axis. Use an integral (in the variable $y$) to find the volume of this solid. 


\newpage



\end{enumerate}

\end{comment} 
\end{enumerate} 

\end{document}
